% !TeX root = ../main.tex
\clearpage
\subsection{تمرین \secnum{۴.۹}، اندازه‌گذاری \lr{NAND4}}
\label{sec:delay-exercise-nand4}

\textbf{خواسته.}
\lr{NAND} چهارورودی را طوری اندازه‌گذاری کنید که مقاومت
مؤثر مسیر صعود و افت آن با \lr{Inverter} واحد برابر باشد.
سپس \(g\) را به دست آورید
\cite[تمرین \secnum{۴.۹}، صفحه ۱۷۷]{weste}.

\textbf{حل.}
در شکل \ref{fig:delay-exercise-nand4}، برای افت خروجی،
چهار \lr{nMOS} سری باید روشن باشند. اگر عرض نسبی هر کدام
۴ باشد، مقاومت کل \(4(R/4)=R\) است. در بدترین حالت صعود،
فقط یکی از چهار \lr{pMOS} موازی روشن است. عرض نسبی ۲
برای هر \lr{pMOS} مقاومت مسیر صعود را به \(2R/2=R\)
می‌رساند.

\begin{figure}[htbp]
  \centering
  \begin{latin}
  \begin{circuitikz}[x=0.94cm,y=0.84cm,font=\small]
    \draw (-0.28,1.02) -- (7.8,1.02);
    \node[above] at (3.9,1.02) {\(V_{DD}\)};
    \draw (0,-1.02) -- (7.8,-1.02) -- (8.55,-1.02)
      node[right] {\(Y\)};
    \foreach \x/\pin in {0/A,2.6/B,5.2/C,7.8/D}{
      \node[pmos] (p\pin) at (\x,0) {};
      \draw (p\pin.S) -- (\x,1.02)
        (p\pin.D) -- (\x,-1.02);
      \fill (\x,1.02) circle (1.25pt)
        (\x,-1.02) circle (1.25pt);
      \node[left] at (p\pin.G) {\(\pin\)};
      \node[right] at (p\pin) {\rl{۲}};
    }
    \foreach \y/\pin in {-2.05/A,-3.25/B,-4.45/C,-5.65/D}{
      \node[nmos] (n\pin) at (3.9,\y) {};
      \node[left] at (n\pin.G) {\(\pin\)};
      \node[right] at (n\pin) {\rl{۴}};
    }
    \draw (3.9,-1.02) -- (nA.D)
      (nA.S) -- (nB.D) (nB.S) -- (nC.D) (nC.S) -- (nD.D)
      (nD.S) -- (3.9,-6.65) node[ground] {};
    \fill (3.9,-1.02) circle (1.35pt);
  \end{circuitikz}
  \end{latin}
  \caption{ساختار ترانزیستوری \lr{NAND4} با چهار \lr{pMOS}
  موازی و چهار \lr{nMOS} سری. حروف \(A\) تا \(D\) پایه‌های
  ورودی‌اند و عدد کنار هر ترانزیستور عرض آن برحسب عرض واحد است.}
  \label{fig:delay-exercise-nand4}
\end{figure}

هر ورودی، \lr{Gate} یک \lr{nMOS} با عرض ۴ و یک \lr{pMOS}
با عرض ۲ را راه‌اندازی می‌کند. پس ظرفیت یک پایه ورودی \(6C\)
است. با مقایسه با \(3C\) برای \lr{Inverter} هم‌قدرت،
\(g=6C/(3C)=6/3=2\) به دست می‌آید.
